% Free asymmetric top, I1 < I2 < I3 (drawn with I1:I2:I3 = 1:2:3).
% Left : exact trajectories of M on the angular-momentum sphere |M| = 1 (intersections with the
%        energy ellipsoid). Visible parts solid, hidden parts dashed. The separatrices through
%        +-x2 are the two great circles n3 = +-sqrt(I3(I2-I1)/(I1(I3-I2))) n1.
% Right: a book-shaped box (a > b > c) and the stability of rotation about its three axes.
% Compile: xelatex asymmetric_top_stability_tikz.tex ; pdftoppm -png -r 300 -singlefile asymmetric_top_stability_tikz.pdf asymmetric_top_stability
\documentclass[tikz,border=4pt]{standalone}
\usepackage{amsmath}
\usepackage[fontset=fandol]{ctex}
\usepackage{pgfplots}
\usepackage{tikz-3dplot}
\pgfplotsset{compat=1.18}
\usetikzlibrary{arrows.meta}

\begin{document}
\begin{tikzpicture}[>=Stealth,line cap=round,every node/.style={font=\normalsize}]
  % ------------------------------------------------------------------ left panel
  \def\az{125}\def\el{22}
  % unit vector towards the viewer for pgfplots view={az}{el}
  \pgfmathsetmacro{\qVX}{sin(\az)*cos(\el)}
  \pgfmathsetmacro{\qVY}{-cos(\az)*cos(\el)}
  \pgfmathsetmacro{\qVZ}{sin(\el)}
  % two unit vectors e, f perpendicular to v (for the sphere's outline)
  \pgfmathsetmacro{\qEX}{cos(\az)}\pgfmathsetmacro{\qEY}{sin(\az)}\pgfmathsetmacro{\qEZ}{0}
  \pgfmathsetmacro{\qFX}{-sin(\az)*sin(\el)}\pgfmathsetmacro{\qFY}{cos(\az)*sin(\el)}\pgfmathsetmacro{\qFZ}{cos(\el)}
  \begin{axis}[name=sph,at={(0,0)},anchor=center,hide axis,view={\az}{\el},
      axis equal image,width=10cm,height=10cm,
      xmin=-1.3,xmax=1.3,ymin=-1.3,ymax=1.3,zmin=-1.3,zmax=1.3,
      clip=false,samples=241,samples y=1,domain=0:360,unbounded coords=jump,
      every axis plot/.append style={line width=0.9pt}]
    % outline of the sphere
    \addplot3[blue!60!black,fill=blue!5,line width=1pt]
      ({cos(x)*\qEX+sin(x)*\qFX},{cos(x)*\qEY+sin(x)*\qFY},{cos(x)*\qEZ+sin(x)*\qFZ});
    % ---- curves around x3 (stable, minimum energy): c = f*(1/I2-1/I3) = f/6
    \foreach \f in {0.3,0.75}{
      \foreach \s in {1,-1}{
        \edef\tmp{
        \noexpand\addplot3[green!45!black,
           restrict expr to domain={\noexpand\qVX*sqrt(\f/6/(2/3))*cos(x)+\noexpand\qVY*sqrt(\f)*sin(x)+\noexpand\qVZ*\s*sqrt(1-\f/4*cos(x)^2-\f*sin(x)^2)}{0:2}]
          ({sqrt(\f/4)*cos(x)},{sqrt(\f)*sin(x)},{\s*sqrt(1-\f/4*cos(x)^2-\f*sin(x)^2)});
        \noexpand\addplot3[green!45!black,densely dashed,line width=0.5pt,
           restrict expr to domain={\noexpand\qVX*sqrt(\f/4)*cos(x)+\noexpand\qVY*sqrt(\f)*sin(x)+\noexpand\qVZ*\s*sqrt(1-\f/4*cos(x)^2-\f*sin(x)^2)}{-2:0}]
          ({sqrt(\f/4)*cos(x)},{sqrt(\f)*sin(x)},{\s*sqrt(1-\f/4*cos(x)^2-\f*sin(x)^2)});
        }\tmp
      }
    }
    % ---- curves around x1 (stable, maximum energy): c' = f*(1/I1-1/I2) = f/2
    \foreach \f in {0.3,0.75}{
      \foreach \s in {1,-1}{
        \edef\tmp{
        \noexpand\addplot3[blue!75!black,
           restrict expr to domain={\noexpand\qVX*\s*sqrt(1-3*\f/4*cos(x)^2-\f*sin(x)^2)+\noexpand\qVY*sqrt(\f)*sin(x)+\noexpand\qVZ*sqrt(3*\f/4)*cos(x)}{0:2}]
          ({\s*sqrt(1-3*\f/4*cos(x)^2-\f*sin(x)^2)},{sqrt(\f)*sin(x)},{sqrt(3*\f/4)*cos(x)});
        \noexpand\addplot3[blue!75!black,densely dashed,line width=0.5pt,
           restrict expr to domain={\noexpand\qVX*\s*sqrt(1-3*\f/4*cos(x)^2-\f*sin(x)^2)+\noexpand\qVY*sqrt(\f)*sin(x)+\noexpand\qVZ*sqrt(3*\f/4)*cos(x)}{-2:0}]
          ({\s*sqrt(1-3*\f/4*cos(x)^2-\f*sin(x)^2)},{sqrt(\f)*sin(x)},{sqrt(3*\f/4)*cos(x)});
        }\tmp
      }
    }
    % ---- separatrices: great circles n = cos(t) (1,0,+-sqrt3)/2 + sin(t) (0,1,0)
    \foreach \s in {1,-1}{
      \edef\tmp{
      \noexpand\addplot3[red,line width=1.4pt,
         restrict expr to domain={\noexpand\qVX*cos(x)/2+\noexpand\qVY*sin(x)+\noexpand\qVZ*\s*sqrt(3)*cos(x)/2}{0:2}]
        ({cos(x)/2},{sin(x)},{\s*sqrt(3)*cos(x)/2});
      \noexpand\addplot3[red,densely dashed,line width=0.6pt,
         restrict expr to domain={\noexpand\qVX*cos(x)/2+\noexpand\qVY*sin(x)+\noexpand\qVZ*\s*sqrt(3)*cos(x)/2}{-2:0}]
        ({cos(x)/2},{sin(x)},{\s*sqrt(3)*cos(x)/2});
      }\tmp
    }
    % ---- axes (outside parts) and poles
    \draw[->,gray!60!black,thick] (axis cs:1,0,0) -- (axis cs:1.45,0,0) node[anchor=north east,black] {$M_1$};
    \draw[->,gray!60!black,thick] (axis cs:0,1,0) -- (axis cs:0,1.45,0) node[anchor=west,black] {$M_2$};
    \draw[->,gray!60!black,thick] (axis cs:0,0,1) -- (axis cs:0,0,1.45) node[anchor=south,black] {$M_3$};
    \fill[blue!75!black] (axis cs:1,0,0) circle (2pt);
    \fill[green!45!black] (axis cs:0,0,1) circle (2pt);
    \node[red,font=\large] at (axis cs:0,1,0) {$\times$};
  \end{axis}
  \node[anchor=north,align=center,font=\small] at (0,-3.0)
    {$\boldsymbol M$ 在刚体系中的轨迹（$I_1:I_2:I_3=1:2:3$）\\
     \textcolor{blue!75!black}{蓝：绕 $x_1$}\quad \textcolor{green!45!black}{绿：绕 $x_3$}\quad
     \textcolor{red}{红：过 $\pm x_2$ 的分离轨道}};

  % ------------------------------------------------------------------ right panel
  \tdplotsetmaincoords{65}{125}
  \begin{scope}[shift={(7.6,0.6)},tdplot_main_coords]
    \def\a{1.55}\def\b{1.05}\def\c{0.22}   % half edges along x1, x2, x3
    % visible faces of the box (top, two front sides)
    \fill[blue!12,draw=blue!60!black] (-\a,-\b,\c) -- (\a,-\b,\c) -- (\a,\b,\c) -- (-\a,\b,\c) -- cycle;
    \fill[blue!22,draw=blue!60!black] (\a,-\b,-\c) -- (\a,\b,-\c) -- (\a,\b,\c) -- (\a,-\b,\c) -- cycle;
    \fill[blue!30,draw=blue!60!black] (-\a,\b,-\c) -- (\a,\b,-\c) -- (\a,\b,\c) -- (-\a,\b,\c) -- cycle;
    % axes
    \draw[->,green!45!black,very thick] (\a,0,0) -- (2.6,0,0) node[anchor=north east] {$x_1$};
    \draw[->,red,very thick] (0,\b,0) -- (0,2.1,0) node[anchor=west] {$x_2$};
    \draw[->,green!45!black,very thick] (0,0,\c) -- (0,0,1.9) node[anchor=south] {$x_3$};
    % rotation arrows about each axis
    \draw[->,green!45!black,thick] plot[domain=30:320,samples=40,variable=\t] ({2.25},{0.3*cos(\t)},{0.3*sin(\t)});
    \draw[->,red,thick] plot[domain=30:320,samples=40,variable=\t] ({0.3*sin(\t)},{1.75},{0.3*cos(\t)});
    \draw[->,green!45!black,thick] plot[domain=30:320,samples=40,variable=\t] ({0.3*cos(\t)},{0.3*sin(\t)},{1.55});
  \end{scope}
  \node[anchor=north,align=left,font=\small] at (7.6,-1.2)
    {长方体（书本、手机）边长 $a>b>c$：\\[2pt]
     \textcolor{green!45!black}{$x_1$（沿最长边）：$I_1$ 最小，稳定}\\
     \textcolor{red}{$x_2$（沿次长边）：$I_2$ 居中，不稳定}\\
     \textcolor{green!45!black}{$x_3$（垂直于大面）：$I_3$ 最大，稳定}};
\end{tikzpicture}
\end{document}
